%% =====================================================================
%%  Temporally Constrained, Provenance-Verified Evidence Retrieval
%%  for Indian Legal Research
%%  IEEE conference manuscript based on the official IEEE conference template.
%% =====================================================================

\documentclass[conference]{IEEEtran}

\usepackage{cite}
\usepackage{amsmath}
\usepackage{graphicx}
\usepackage{booktabs}
\usepackage{xurl}   % allow long URLs to break anywhere (repo link in Sec. Reproducibility)
\usepackage[hidelinks]{hyperref}

\def\BibTeX{{\rm B\kern-.05em{\sc i\kern-.025em b}\kern-.08em
    T\kern-.1667em\lower.7ex\hbox{E}\kern-.125emX}}

% Tighter tables in two-column layout
\renewcommand{\arraystretch}{1.12}

\begin{document}

\title{Temporally Constrained, Provenance-Verified Evidence Retrieval for Indian Legal Research}

\author{
\begin{tabular}{c@{\hspace{1.0em}}c}
Yuvraj Singh & RishiRaj Jaiswal \\
\textit{Department of Computer Science and Engineering} & \textit{Department of Computer Science and Engineering} \\
\textit{Ajay Kumar Garg Engineering College} & \textit{Ajay Kumar Garg Engineering College} \\
Ghaziabad, India & Ghaziabad, India \\
{\small yuvraj2312026@akgec.ac.in} & {\small rishi2312003@akgec.ac.in} \\[2ex]
Rahul Ranjan & Saurav Kumar Chaudhary \\
\textit{Department of Computer Science and Engineering} & \textit{Department of Computer Science and Engineering} \\
\textit{Ajay Kumar Garg Engineering College} & \textit{Ajay Kumar Garg Engineering College} \\
Ghaziabad, India & Ghaziabad, India \\
{\small rahul2312151@akgec.ac.in} & {\small saurav2312023@akgec.ac.in}
\end{tabular}
}

\maketitle

% =====================================================================
\begin{abstract}
AI-assisted legal research fails in ways that fluent output conceals: a cited
authority may not exist, may lack the attributed passage, may post-date
the matter, or may never have been retrieved. A 2026 Supreme Court of India
ruling (2026 INSC 668) overturned tribunal findings built on exactly these
defects. We present an
evidence pipeline for
historical Indian Supreme Court judgments in which every displayed citation is
temporally eligible, provenance-bound, and mechanically verifiable, with outcome
prediction kept separate. We report three independently replayed findings. First, naive ILDC--eCourts identifier matching proves unsafe, as only 11 out of 5{,}391 syntactic matches survived strict content alignment.
Second, pre-ranking temporal eligibility
raised Recall@5 from
5/30 to 12/30 and Recall@100 from 12/30 to 15/30 on the frozen 30-case base.
Third, all 150 displayed base citations passed verification,
with zero unsupported claims,
while 15 expected authorities remained
absent at $k{=}100$. Outcome baselines on 1{,}503 cases use their own population.

\end{abstract}

\begin{IEEEkeywords}
legal information retrieval, explainable AI, provenance verification, temporal
filtering, retrieval-augmented systems, Indian legal NLP, BM25, InLegalBERT
\end{IEEEkeywords}

% =====================================================================
\section{Introduction}

Legal research rewards fluency over evidence: generated
responses cite nonexistent judgments, misattribute passages, or
invoke unavailable law. Such hallucinations are frequent, not exceptional
\cite{dahl2024}, and communicating legal standards to models remains open
\cite{nay2023}.

The problem is no longer confined to the literature. In \emph{Pooja Ramesh Singh
v.\ Jammu and Kashmir Bank Ltd.\ \& Anr.}, 2026 INSC 668 (Civil Appeal
No.\ 11950 of 2025, decided 2 July 2026), the Supreme Court of
India set aside tribunal decisions after identifying nonexistent authorities,
incorrect citations, and fabricated passages apparently produced through
AI-assisted research \cite{poojasingh2026}.

We treat that constraint as the design foundation:
retrieved facts carry stable
locators, duplicates are excluded, only verbatim passages are rendered,
and each citation is verified against corpus and retrieval run.

Our primary finding highlights a stark contrast: while all 150 displayed citations passed strict verification without a single unsupported claim, 15 of the 30 expected authorities failed to appear within the top 100 retrieved results.

\textbf{Contributions.} We make four main contributions:

\begin{enumerate}
\item \textbf{A quantified cross-corpus identity failure and its remedy.}
\texttt{YYYY INSC N} to \texttt{YYYY\_N} conversion produced 5{,}391 syntactic candidates, of which only 11 passed
content alignment (Section~\ref{sec:alignment}).

\item \textbf{Pre-ranking temporal eligibility.} Historical
availability as a predicate inside the BM25 candidate relation, so ineligible judgments cannot consume retrieval depth
(Section~\ref{sec:e3}): Recall@5 moved
from 5/30 to 12/30 without regressing any prior success.

\item \textbf{A four-state citation protocol with separate denominators.}
Fabricated, real-but-unretrieved, policy-violating, and valid-but-unexpected
citations (Section~\ref{sec:protocol}).

\item \textbf{An evaluation that keeps recovery and integrity apart.} Authority
recovery, citation integrity, outcome prediction, and
presentation are reported on their own populations
(Section~\ref{sec:results}).
\end{enumerate}

\textbf{Scope.} No legal
correctness for retrieved authorities, no autonomous legal advisor, no
claim that retrieval improves outcome prediction: a
researcher-facing prototype on four non-pooled
populations (1{,}503 outcome cases; 30 frozen evidence cases comprising the replayed core).

% =====================================================================
\section{Related Work}

\subsection{Indian legal NLP and judgment prediction}

ILDC established Indian Supreme Court judgment prediction and
explanation as a benchmark task \cite{malik2021}; InLegalBERT \cite{paul2023} motivates our neural
baseline; practitioner-evaluated summarisation \cite{shukla2022}
shows length and presentation are first-order concerns. None of
these lines addresses cross-corpus identity or per-citation
provenance, which is where this work sits.

\subsection{Legal retrieval and temporal constraints}

COLIEE, LegalBench, and CaseHOLD evaluate retrieval and reasoning on single
fixed corpora with no pre-dating requirement \cite{coliee2023,legalbench2023,casehold2021}.
TaxFlow adds validity filtering for Indian tax-law QA \cite{taxflow2026};
CaseFacts and LexTime target evolving validity and event ordering
\cite{casefacts2026,lextime2025}. Retrieval-augmented generation
established retrieval as an explicit non-parametric evidence source
\cite{lewis2020}; BM25 supplies the sparse ranking basis \cite{robertson2009}
and Tesseract the OCR engine for corpus repair \cite{smith2007}.

\subsection{Positioning}

We claim neither the first legal RAG system, nor the first
Indian legal model, nor the first temporal legal benchmark: the contribution is
the reproducible combination of content-alignment gating,
conservative pre-ranking eligibility, exact retrieval-run
verification, and separate recovery-versus-validity reporting on a source-first key.
Jurisdictions, tasks, corpora, and denominators differ, so no external system is
used as a numerical baseline. ILDC \cite{malik2021} supplies splits; InLegalBERT \cite{paul2023} motivates the neural baseline;
TaxFlow \cite{taxflow2026} addresses statutory validity; CaseFacts \cite{casefacts2026} and LexTime
\cite{lextime2025} target validity and event ordering.

% =====================================================================
\section{Problem Formulation}
\label{sec:problem}

Let a historical query judgment have identifier $q$, query year $Y_q$, and
frozen facts-only text $F_q$. Let each candidate evidence passage $p$ belong to
a source judgment with a stable source identifier, an exact decision date,
citation metadata, page and character boundaries, and stored passage text.

The temporal policy admits $p$ only when the decision year of its source is
strictly earlier than $Y_q$. A same-year source is treated as ambiguous, because
ILDC supplies no exact query date; a source with missing or unparseable date
metadata is excluded. Both rules trade
recall so that, under the fail-closed metadata checks applied here, no
displayed authority can post-date the matter.

The evidence task decomposes into recovery (expected authority in the top 100
and among the five displayed?),
integrity (each citation a real retrieved passage with
exact provenance, avoiding query/duplicates, under the temporal rule?),
and presentation (does structure aid inspection?).

Outcome prediction is secondary, reported on its own
population and never pooled with evidence measures. Table~\ref{tab:defs} fixes
definitions.

\subsection{Research questions}

We investigate two primary research questions. \textbf{RQ1:} Does grounding a
legal AI workflow in retrieved evidence improve reliability and retrieval against a facts-only
baseline (E2 versus E3; E1 for context)? \textbf{RQ2:} Does
provenance-constrained selection and verification reduce
unsupported claims?


\begin{table}[t]
\caption{Operational definitions (frozen evaluation).}
\label{tab:defs}
\centering
\footnotesize
\begin{tabular}{p{2.35cm}p{5.5cm}}
\toprule
\textbf{Term} & \textbf{Implemented definition} \\
\midrule
Facts-only input &
Text before the earlier of a recognised dispositive cue and a 60\%
character cap; inputs below 10\% retention or
100 words are excluded \\
\addlinespace
Temporal existence &
The source has a parseable exact decision date \\
\addlinespace
Temporal applicability &
For query year $Y$, an authority is eligible only if its decision year is
$< Y$; same-year, later-year, and missing-date sources are excluded \\
\addlinespace
Temporal effectiveness &
Eligibility is applied inside the BM25 candidate relation before
ranking and \texttt{LIMIT 100} \\
\addlinespace
Recall@100 &
Share of cases whose predefined authority occurs in the
returned top-100 \\
\addlinespace
Recall@5 &
Share of cases whose predefined authority is among the five selected
sources \\
\addlinespace
Provenance validity &
A displayed item reproduces the stored source ID, citation, decision date,
court, locator, exact passage, and retrieval-run membership \\
\addlinespace
Citation groundedness &
Every displayed material proposition is verbatim text of a supplied passage
linked to that evidence item \\
\addlinespace
Authority consistency &
A displayed source matches the answer-key authority by stable source
ID, normalised citation, or title plus exact decision date \\
\bottomrule
\end{tabular}
\end{table}

% =====================================================================
\section{Corpus Construction and Auditing}

\subsection{Purpose-split corpus design}

We use two corpora for different experimental purposes.

ILDC Single supplies fixed case-level splits and binary outcome labels: 7{,}593 judgments (5{,}082 training, 994 validation,
1{,}517 test). After a shared sufficiency rule excluded 14 test records,
the prediction population contains 1{,}503 cases.

The evidence corpus is an eCourts-derived collection of Indian Supreme Court
judgments covering 1950--2020. Of 39{,}069 English
PDF instances, 39{,}066 accepted PDFs yielded
2{,}343{,}435 labelled chunks, of which 2{,}036{,}981 unique chunks were loaded
into a PostgreSQL provenance store and a SQLite FTS5 BM25 index.

This split follows the information actually available in each source. ILDC
supports fixed-split outcome prediction but lacks citation and passage
provenance; the eCourts collection supports dated, traceable retrieval but is
not an outcome-label benchmark. Cross-corpus links are drawn only
after the alignment checks described below.

\subsection{Source-quality audit and targeted OCR repair}

A full audit found valid UTF-8 text except 15 PDF instances
(14 source IDs) with missing or corrupted embedded text.

Only flagged instances were re-rendered and processed with English Tesseract OCR
at 250\,DPI \cite{smith2007}. Instances passing a post-OCR quality gate were restored. Raw PDFs were preserved and the index rebuilt from the accepted corpus.

\subsection{Cross-corpus identifier alignment}
\label{sec:alignment}

Linking ILDC query cases to eCourts sources requires establishing that two
records describe the same judgment. The obvious route --- converting an eCourts
identifier of the form \texttt{YYYY INSC N} to the ILDC-like string
\texttt{YYYY\_N} --- is unsafe, and quantifiably so.

That conversion produced 5{,}391 syntactic candidates. \textbf{Only 11 passed
content alignment; 5{,}380 were identifier-namespace collisions.}

The corrected pipeline treats identifier equality as candidate generation only.
A reusable gate evaluates title/party identity with direct six-token
phrase overlap before any mapping drives deduplication or
exclusion. Combining syntactic and title/party candidates yielded 8{,}927
deduplicated pairs (1{,}304 accepted, 7{,}623 rejected).
Retrieval additionally applies a full-document self-match safeguard, so an
unmapped copy of the query judgment is removed without suppressing an
earlier authority quoted at length by the later judgment.

\subsection{Source-verified authority answer key}

The authority evaluation uses 30 ILDC fixed-test cases. For each, the query
judgment was inspected through a primary source or accepted eCourts mirror, one
earlier authority was recorded \emph{independently of system retrieval}, and
reconciled by stable source ID, normalised
citation, or title plus exact date.

The identifier-collision discovery triggered a read-only audit of the
then-current mappings: 20 passed, nine failed direct content
alignment, and one was unresolved. One case was relinked
and the other nine replaced by new
fixed-test cases, yielding 30/30 direct-content passes.

The sample is not era-balanced: five query cases are from the 1970s, 13 from the
1980s, nine from the 1990s, two from the 2000s, and one from the 2010s
(1980s: 43.3\%), a consequence of the source and
alignment gates (Section~\ref{sec:limits}).


% =====================================================================
\section{System Design}

Four auditable layers implement the design. A corpus and alignment layer links
ILDC cases to dated eCourts passages with content-checked deduplication. An
outcome branch (E1 sparse linear, E2 InLegalBERT chunk-and-pool) predicts from
facts only and never retrieves. An evidence branch (E3) retrieves BM25 passages
under pre-ranking temporal eligibility and renders five source-diverse verbatim
passages. A verification layer (E4) checks each displayed record against the
corpus and the recorded retrieval run. The renderer is extractive and cannot
retrieve, paraphrase material text, or infer outcomes; prediction stays separate.

% =====================================================================
\section{Method}

\subsection{Shared facts-only input}

E1 and E2 share facts-only extraction (\texttt{ildc-predecision-facts-v1}): text preceding a dispositive cue
and 60\% of the document. Cases below 10\% retention or 100 words are excluded;
selection uses
validation only, with one held-out test evaluation.

\subsection{E1: sparse linear baseline}

E1 uses lowercased TF--IDF unigrams/bigrams (sublinear TF, min DF two, $L_2$,
100{,}000 features). Of
$C \in \{0.1, 1.0, 10.0\}$, validation accuracy selected $C = 10.0$ (smaller $C$
ties); refit once on eligible
training plus validation and evaluated once on test. Seed 202605.

\subsection{E2: InLegalBERT chunk-and-pool}

E2 fine-tunes \texttt{law-ai/InLegalBERT} (revision
\texttt{b5ecfed8}) with a new two-label head over overlapping 512-token windows with 50-token overlap.
Training uses three epochs, seed 202607, learning rate $2\times10^{-5}$, weight
decay 0.01, warm-up ratio 0.1, and gradient
accumulation.

Mean pooling of window logits before softmax is the primary rule, and checkpoint
selection uses validation document-level mean-logit accuracy only (majority vote
is a secondary comparison, not a selection rule). All
1{,}503 eligible test documents and 9{,}576 test windows are covered.

\subsection{E3: temporally constrained retrieval}
\label{sec:e3}

The query builder (\texttt{tfidf-\allowbreak segment-\allowbreak salient-\allowbreak terms-\allowbreak v1}) emits at most 32 salient
terms for a SQLite FTS5 BM25 index backed by PostgreSQL provenance.

The frozen configuration (\texttt{week11-\allowbreak bm25-\allowbreak salient-\allowbreak terms-\allowbreak preranked-\allowbreak temporal-\allowbreak v3}) admits only judgments with decision year strictly earlier than the
query year into the candidate relation \emph{before} BM25 ordering and the
top-100 cutoff, so the full depth is spent on eligible material.

Candidate sources are checked against the alignment-gated crosswalk and a
self-match rule (100 shared six-token occurrences, 80\%
source-phrase coverage). A non-learned selector then chooses up
to five passages in BM25 order, at most one per source.

\subsection{E4: citation and provenance verification}

E4 receives the E3 answer, retrieval-run identifier, and
persisted corpus records. Each displayed item must exist in the
corpus, reproduce exact passage and provenance, belong to the recorded
run, avoid the query and duplicates, and pre-date the
query year; any failure rejects the citation. Authority consistency is then
evaluated separately: a fully valid citation may differ from the single
expected authority without being substantively irrelevant (Section~\ref{sec:ceiling}).

\subsection{Inference-time evidence augmentation}

Outcome prediction was subsequently added to E3 and E4 by reusing the frozen E2
checkpoint without fine-tuning: the frozen
facts extract plus the selected verbatim evidence in selection order pass through
the same 512-token/50-overlap windows with mean-logit \texttt{argmax} (seed
202607; configuration
\texttt{e3e4-\allowbreak evidence-\allowbreak augmented-\allowbreak inlegalbert-\allowbreak checkpoint6318-\allowbreak v1}).
Gold labels, scores, dates, and provenance are excluded from model input.

We note the caveat plainly: the checkpoint was trained on
facts-only inputs, so applying it to facts plus retrieved passages is an
inference-time distribution shift, not evidence-aware fine-tuning. These
predictions are reported descriptively on the frozen 30-case
base, and are
not directly comparable with the 1{,}503-case baselines.

\subsection{Bundled-intervention caveat}

E4 jointly checks corpus existence,
passage/provenance identity, run membership, duplicate
status, and temporal eligibility. Aggregate reliability cannot be attributed to any one
component; the supported claim concerns the complete bundle.

% =====================================================================
\section{Citation and Provenance Protocol}
\label{sec:protocol}

Each chunk carries a stable identifier from its source path and
passage location, with stored source ID, case ID, citation, exact decision date, court, PDF
file, page and character bounds, and chunk text. Retrieval runs record query ID, year,
query text, index version, temporal policy, rank, score, and temporal status.

The renderer may expose only supplied selected evidence. Material content
remains verbatim; authority cards and conclusions refer to evidence IDs.
Verification then checks corpus existence, exact field equality, run membership,
duplicate status, and temporal eligibility, with answer-key matching only
after a displayed citation has passed. Parallel reporter forms are
reconciled through stable source identity or normalised title plus exact
decision date, never loose string similarity.

This protocol distinguishes four failure states that a single
``correct citation'' label would collapse:

\begin{enumerate}
\item a fabricated or altered citation;
\item a real citation that was not retrieved for this query;
\item a retrieved citation violating temporal or duplicate policy;
\item a fully valid citation that does not match the predefined expected
authority.
\end{enumerate}

In our results, states (1)--(3) occur zero times and state (4) occurs
in 18 of 30 cases.

% =====================================================================
\section{Results}
\label{sec:results}

\subsection{Reporting populations}

Two populations are reported and never pooled: 1{,}503
ILDC test cases (outcome prediction) and Base-30 (30 queries; 150 citations) for evidence evaluation.

\subsection{Outcome prediction}

On the 1{,}503-case population (Table~\ref{tab:outcome}), sparse E1 reached accuracy 0.61344 and macro-F1 0.612342,
exceeding corrected E2 mean-logit pooling at 0.596806 and 0.592358 (majority vote
0.6015 and 0.5937). All primary models exceeded
the 0.5017 majority-class accuracy.

\begin{table}[t]
\caption{Outcome prediction (distinct populations; not a leaderboard).}
\label{tab:outcome}
\centering
\footnotesize
\setlength{\tabcolsep}{4pt}
\begin{tabular}{lccc}
\toprule
\textbf{Model} & \textbf{$n$} & \textbf{Acc.} & \textbf{Macro-F1} \\
\midrule
Majority class            & 1{,}503 & 0.5017 & 0.3341 \\
E1 TF--IDF + LR           & 1{,}503 & \textbf{0.61344} & \textbf{0.612342} \\
E2 InLegalBERT, mean logits & 1{,}503 & 0.596806 & 0.592358 \\
E2 InLegalBERT, maj.\ vote  & 1{,}503 & 0.6015 & 0.5937 \\
\midrule
E3 evidence-augmented     & 30 & 0.666667 & 0.603175 \\
E4 verified evid.-augmented & 30 & 0.666667 & 0.603175 \\
\bottomrule
\end{tabular}
\end{table}


While this does not imply that domain pre-training is generally ineffective, it does show that the sparse baseline performed better under our frozen settings.

On the separate 30-case subset, E3 and E4 both
reached accuracy 0.666667 and macro-F1 0.603175, with identical confusion matrices.

\textbf{The equality is expected:} with no evidence rejected, the predictor received an identical
input.

\subsection{Authority recovery and evidence integrity}

Base-30 recovery is Recall@5 0.40 (12/30) and Recall@100 0.50 (15/30: 12
selected $+$ 3 retrieved-but-unselected); 15 authorities were absent at
$k{=}100$.

The principal RQ1 result is the Base-30 row of Table~\ref{tab:rq1strata}.


\begin{table}[t]
\caption{RQ1 retrieval and prediction by stratum (Base-30 replayed).}
\label{tab:rq1strata}
\centering
\footnotesize
\setlength{\tabcolsep}{2pt}
\begin{tabular}{lrrrrr}
\toprule
\textbf{Population} & \textbf{$n$} & \textbf{R@5} & \textbf{R@100} &
\textbf{E3 Acc.} & \textbf{E3 F1} \\
\midrule
Base-30          & 30 & 0.40 (12/30) & 0.50 (15/30)   & 0.6667 & 0.6032 \\

\bottomrule
\end{tabular}
\end{table}

Table~\ref{tab:integrity} summarizes Base-30 integrity:
150/150 grounded and provenance-valid with 0/150 temporal violations and 0/30
unsupported claims.


E4 verification is fail-closed: any failing check rejects the citation. On
live data nothing was rejected, a ceiling effect of already-valid evidence.

\begin{table}[t]
\caption{Evidence recovery and integrity, Base-30 (30
queries; 150 citations).}
\label{tab:integrity}
\centering
\footnotesize
\setlength{\tabcolsep}{4pt}
\begin{tabular}{lcl}
\toprule
\textbf{Measure} & \textbf{Value} & \textbf{Arithmetic} \\
\midrule
Recall@5            & 0.40 & 12 of 30 in top 5 \\
Recall@100          & 0.50 & 15 of 30 in top 100 \\
Absent at $k{=}100$ & 0.50 & 15 of 30 \\
Authority precision & 0.08 & 12 of 150 shown \\
\quad struct.\ ceiling & 0.20 & 30 of 150 (5 shown, 1 credited) \\
Grounding           & 1.00 & 150/150 passed \\
Provenance validity & 1.00 & 150/150 passed \\
Temporal violations & 0    & 0 of 150 \\
Unsupported claims  & 0    & 0 of 30 cases \\
\bottomrule
\end{tabular}
\end{table}



\textbf{This is the paper's central result.} The system verified everything it displayed (150/150)
while recovering the expected authority
in 15 of 30 cases at $k{=}100$.

\subsection{Interpreting authority-consistency precision}
\label{sec:ceiling}

Authority-consistency precision was 0.08, with recall 0.40 and F1 0.133:
with one credited authority per case and five displayed, perfect inclusion caps
the measure at $30/150 = 0.20$, so 0.08 is 40\% of ceiling.

The 138 non-matching displayed citations are \emph{not} independently annotated
for substantive legal relevance. They passed
every integrity check; they simply are not the one authority the key records.

\subsection{Retrieval mechanism}

The final configuration came from three recorded development steps on a
nine-case probe (0/9 to 3/9 to 6/9 to 7/9 at $k{=}100$); pre-ranking temporal filtering,
carried to the held-out base, moved Recall@5 5/30 to 12/30 and Recall@100 12/30 to
15/30 \emph{without losing any of the 12 prior top-100 successes}. Probe history is
reported from development records, not replayed.


Temporal capacity is not the whole story. \texttt{2013\_35} stayed absent
despite 79 eligible candidates in its raw top 100, while
\texttt{1980\_105} was recovered from only three and
became a top-five hit: residual lexical mismatch is the prominent unresolved
cause.

% =====================================================================
\section{Error Analysis}
\label{sec:error}

\subsection{Outcome disagreement}

The analysis is read-only over final frozen outputs. Across 1{,}503 cases, both models were correct on 684
and both wrong on 368; E1 alone was correct on 238 and E2 alone on 213
(on the 30-case key: 18/3/3/6).


In the evidence-augmented run, E2 was wrong while E3 and E4 were correct for two
cases (\texttt{1974\_36}, \texttt{1984\_136}); E3/E4 verification changed
neither the selected evidence nor the shared predictor input (0/30 divergences).

Critically, three cases (\texttt{2008\_1629}, \texttt{1981\_187},
\texttt{1982\_29}) retrieved \emph{and} selected the expected
authority yet still produced an incorrect outcome prediction, and a fourth
(\texttt{1985\_40}) was retrieved but not selected yet likewise predicted
incorrectly: correct recovery does not mechanically imply a correct prediction. No correct prediction was
accompanied by an unsupported citation.

\subsection{Authority-recovery buckets}

The exhaustive partition is: 12 retrieved and
selected; 3 retrieved but unselected
(\texttt{1980\_133} rank 15, \texttt{1981\_55} rank 28,
\texttt{1985\_40} rank 78); and 15 absent at $k{=}100$.
The middle bucket isolates \emph{selection} failure (available but undisplayed), while the absent bucket requires upstream
improvement.


% Worked-example detail retained in frozen artifacts; subsection removed for length.

\section{Limitations and Threats to Validity}
\label{sec:limits}

\textbf{Scope.} English-language Indian Supreme Court judgments
only.

\textbf{Extraction approximation.} Some affected PDF instances were restored through the post-OCR
quality gate \cite{smith2007}; repaired text may retain noise. With
no gold facts/reasoning boundary, the facts rule is approximate.

\textbf{Answer-key size and era concentration.} Evidence evaluation rests on
30 frozen replayed cases; 13 of the base 30 are from
the 1980s. The evaluated set is neither large nor era-balanced and should not
be read as representative of the full ILDC test set or of
Indian legal research generally.

\textbf{Single reference authority.} One expected authority per
query, not every relevant source: inconsistency is a
reference mismatch, not a legal-relevance judgment.

\textbf{Residual recovery gap.} At $k{=}100$, 15 of the 30 expected base authorities
remain unrecovered; the \texttt{2013\_35}/\texttt{1980\_105} contrast shows
temporal room alone does not surface the authority.

\textbf{Year-level granularity.} Query dates have
year-level precision only, so same-year candidates
are conservatively excluded.

\textbf{Bundled verification.} E4 checks integrity properties together
and cannot isolate single-component contributions.

\textbf{Inference numerics.} Recomputed E3/E4 mean logits show bounded GPU
fp16 variation across runs (bounded at the fourth decimal in the
independent recomputation) that changed no decision, evidence selection, or
metric; strict serialized file hashes of E3/E4 outputs therefore differ while
substantive outputs agree.

\textbf{Retrieval-equivalence scope.} BM25 rebuild equivalence is established
at ranking level --- identical corpus, parameters, and top-100 agreement on
all 30 evaluated queries --- not as byte-identical SQLite files, and it does
not generalize beyond the evaluated query set.



\textbf{Prototype scale.} Results apply to the frozen implementation and samples.
No production deployment is claimed or evaluated.

% =====================================================================
\section{Responsible Use and Governance}

The system is a legal-research aid, not an adjudicator or provider
of legal advice. E3/E4 predictions are experimental; a
human reviewer judges relevance, currency, authority, and
applicability \cite{poojasingh2026}. Controls are fail-closed and negative results
(discarded runs, collision correction, key replacements) are preserved.

% =====================================================================
\section{Conclusion and Future Work}

We designed, implemented, and evaluated a provenance-preserving, temporally constrained evidence
workflow for historical Indian Supreme Court research.

The principal finding is a dissociation. Every one of 150 displayed base citations
passed verification, with zero unsupported claims,
while 15 of
the 30 expected authorities were
absent at $k{=}100$.

The cross-corpus audit yields a transferable lesson: only 11 of
5{,}391 syntactic identifier matches between ILDC and the eCourts collection
were content-aligned.

We report supporting results along with their limits. The sparse TF--IDF baseline outperformed the corrected InLegalBERT under the frozen settings (1{,}503 cases). Inference-time evidence augmentation
reached 0.666667 accuracy and 0.603175 macro-F1 on the frozen 30-case subset,
descriptively and under acknowledged distribution shift.

\textbf{Future work.} An era-balanced key of 75--100 cases for paired testing;
blinded review with independent raters; hybrid
dense--sparse retrieval for residual lexical misses; and evidence-sensitive
uncertainty language.

% =====================================================================
\section*{Reproducibility}

All configurations, seeds, model revisions, corpus identities, and evaluation
artifacts are frozen and machine-readable, including checkpoint-6318
(\texttt{924a5bb9\ldots dbdc773}). The implementation passes 81
tests and the documented
checks (14/14 documentation; freeze audit of 39 entries:
28 byte-exact, one line-ending-normalized, 10 metadata-differing).
All independently replayed primary metrics matched their frozen references; E3/E4
decisions, evidence selections, and metrics matched despite bounded fp16 variation,
and BM25 matched at ranking level on the 30 evaluated queries.
The demo stack is covered by 8/8 Spring Boot service tests and 6/6 FastAPI
wrapper tests. Code and artifacts are available
at \url{https://github.com/0-YuvrajSingh/NyayaTrace}.

% =====================================================================
\section*{Acknowledgment}

The authors thank Kriti Mishra for her mentorship and guidance throughout this work.

\begin{thebibliography}{00}
\bibitem{dahl2024}
M. Dahl, V. Magesh, M. Suzgun, and D. E. Ho, ``Large Legal Fictions: Profiling
Legal Hallucinations in Large Language Models,'' \emph{Journal of Legal
Analysis}, vol.\ 16, no.\ 1, pp.\ 64--93, 2024.
\bibitem{nay2023}
J. J. Nay, ``Large Language Models as Fiduciaries: A Case Study Toward Robustly
Communicating With Artificial Intelligence Through Legal Standards,''
\emph{arXiv preprint arXiv:2301.10095}, 2023.
\bibitem{poojasingh2026}
Supreme Court of India, \emph{Pooja Ramesh Singh v. Jammu and Kashmir Bank Ltd.\
\& Anr.}, 2026 INSC 668 (Civil Appeal No.\ 11950 of 2025, decided 2 July 2026),
2026. Court PDF:
\url{https://www.sci.gov.in/sci-get-pdf/?diary_no=523382025&from=latest_judgements_order&order_date=2026-07-02&type=j}.
\bibitem{malik2021}
V. Malik, R. Sanjay, S. K. Nigam, K. Ghosh, S. K. Guha, A. Bhattacharya, and
A. Modi, ``ILDC for CJPE: Indian Legal Documents Corpus for Court Judgment
Prediction and Explanation,'' in \emph{Proc.\ 59th Annual Meeting of the
Association for Computational Linguistics (ACL-IJCNLP)}, 2021, pp.\ 4046--4062.
\bibitem{paul2023}
S. Paul, A. Mandal, P. Goyal, and S. Ghosh, ``Pre-trained Language Models for
the Legal Domain: A Case Study on Indian Law,'' in \emph{Proc.\ 19th Int.\
Conf.\ on Artificial Intelligence and Law (ICAIL)}, 2023, pp.\ 187--196.
\bibitem{shukla2022}
A. Shukla, P. Bhattacharya, S. Poddar, R. Mukherjee, K. Ghosh, P. Goyal, and
S. Ghosh, ``Legal Case Document Summarization: Extractive and Abstractive
Methods and their Evaluation,'' in \emph{Proc.\ 2nd Conf.\ of the Asia-Pacific
Chapter of the ACL (AACL-IJCNLP)}, 2022, pp.\ 1048--1064.
\bibitem{coliee2023}
R. Goebel, Y. Kano, M.-Y. Kim, J. Rabelo, K. Satoh, and M. Yoshioka, ``Summary
of the Competition on Legal Information, Extraction/Entailment (COLIEE) 2023,''
in \emph{Proc.\ 19th Int.\ Conf.\ on Artificial Intelligence and Law (ICAIL)},
2023, pp.\ 472--480.
\bibitem{legalbench2023}
N. Guha et al., ``LegalBench: A Collaboratively Built Benchmark for Measuring
Legal Reasoning in Large Language Models,'' in \emph{Advances in Neural
Information Processing Systems 36, Datasets and Benchmarks Track}, 2023.
\bibitem{casehold2021}
L. Zheng, N. Guha, B. R. Anderson, P. Henderson, and D. E. Ho, ``When Does
Pretraining Help? Assessing Self-Supervised Learning for Law and the CaseHOLD
Dataset of 53,000+ Legal Holdings,'' in \emph{Proc.\ 18th Int.\ Conf.\ on
Artificial Intelligence and Law (ICAIL)}, 2021, pp.\ 159--168.
\bibitem{taxflow2026}
V. R. Karna, R. R. M., B. S. Babu, \emph{et al.},
``A Hybrid RAG-LLaMA Framework for Scalable and Accurate Interpretation of Legal
Texts,'' \emph{Applied Artificial Intelligence}, vol.\ 40, no.\ 1, art.\
2626097, 2026. DOI: \url{https://doi.org/10.1080/08839514.2026.2626097}.
\bibitem{casefacts2026}
A. R. Putta, J. Devasier, and C. Li, ``CaseFacts: A Benchmark for Legal
Fact-Checking and Precedent Retrieval,'' in \emph{Proc.\ 64th Annual Meeting of
the Association for Computational Linguistics}, 2026, pp.\ 17246--17265.
\bibitem{lextime2025}
C. Barale, L. Barrett, V. S. Bajaj, and M. Rovatsos, ``LexTime: A Benchmark for
Temporal Ordering of Legal Events,'' in \emph{Findings of the Association for
Computational Linguistics: EMNLP 2025}, 2025, pp.\ 5220--5236.
\bibitem{lewis2020}
P. Lewis et al., ``Retrieval-Augmented Generation for Knowledge-Intensive NLP
Tasks,'' in \emph{Advances in Neural Information Processing Systems 33}, 2020,
pp.\ 9459--9474.
\bibitem{robertson2009}
S. Robertson and H. Zaragoza, ``The Probabilistic Relevance Framework: BM25 and
Beyond,'' \emph{Foundations and Trends in Information Retrieval}, vol.\ 3, no.\
4, pp.\ 333--389, 2009.
\bibitem{smith2007}
R. Smith, ``An Overview of the Tesseract OCR Engine,'' in \emph{Proc.\ 9th Int.\
Conf.\ on Document Analysis and Recognition (ICDAR)}, vol.\ 2, 2007, pp.\
629--633.
\end{thebibliography}

\end{document}
